% !TEX root = long.tex
\documentclass[a4paper,11pt]{article}
% Shared layout for marti_ferran_gonzales_cv.tex and marti_ferran_gonzales_cover_letter.tex.
% The CV defines \cvsetupfont before % Shared layout for marti_ferran_gonzales_cv.tex and marti_ferran_gonzales_cover_letter.tex.
% The CV defines \cvsetupfont before % Shared layout for marti_ferran_gonzales_cv.tex and marti_ferran_gonzales_cover_letter.tex.
% The CV defines \cvsetupfont before \input{preamble}.
% Author : Jake Gutierrez
% Based off of: https://github.com/sb2nov/resume
% License : MIT
% This file is not a document. Compile marti_ferran_gonzales_cv.tex or marti_ferran_gonzales_cover_letter.tex.

\ifdefined\cvsetupfont
  \usepackage{fontspec}
  \cvsetupfont
  \renewcommand{\familydefault}{\sfdefault}
\else
  \usepackage[T1]{fontenc}
  \usepackage[utf8]{inputenc}
\fi
\usepackage{latexsym}
\usepackage[empty]{fullpage}
\usepackage{titlesec}
\usepackage{marvosym}
\usepackage[usenames,dvipsnames]{color}
\usepackage{verbatim}
\usepackage{enumitem}
\usepackage[hidelinks]{hyperref}
\usepackage{fancyhdr}
\usepackage[english]{babel}
\usepackage{tabularx}
\ifdefined\cvsetupfont
\else
  \input{glyphtounicode}
\fi

\pagestyle{fancy}
\fancyhf{} % clear all header and footer fields
\fancyfoot{}
\renewcommand{\headrulewidth}{0pt}
\renewcommand{\footrulewidth}{0pt}

% Adjust margins
\addtolength{\oddsidemargin}{-0.5in}
\addtolength{\evensidemargin}{-0.5in}
\addtolength{\textwidth}{1in}
\addtolength{\topmargin}{-.5in}
\addtolength{\textheight}{1.0in}

\urlstyle{same}

\raggedbottom
\raggedright
\setlength{\tabcolsep}{0in}

% Sans faces here have no small caps, so the section title stays in the same font.
\ifdefined\cvsetupfont
  \titleformat{\section}{
    \vspace{-4pt}\raggedright\large
  }{}{0em}{\MakeUppercase}[\color{black}\titlerule \vspace{-5pt}]
\else
  \titleformat{\section}{
    \vspace{-4pt}\scshape\raggedright\large
  }{}{0em}{}[\color{black}\titlerule \vspace{-5pt}]
  % Ensure that the generated pdf is machine readable/ATS parsable
  \pdfgentounicode=1
\fi

%-------------------------
% Custom commands
\newcommand{\resumeItem}[1]{
  \item\small{
    {#1 \vspace{-2pt}}
  }
}

\newcommand{\resumeSubheading}[4]{
  \vspace{-2pt}\item
    \begin{tabular*}{0.97\textwidth}[t]{l@{\extracolsep{\fill}}r}
      \textbf{#1} & #2 \\
      \textit{\small#3} & \textit{\small #4} \\
    \end{tabular*}\vspace{-7pt}
}

\newcommand{\resumeSubSubheading}[2]{
    \item
    \begin{tabular*}{0.97\textwidth}{l@{\extracolsep{\fill}}r}
      \textit{\small#1} & \textit{\small #2} \\
    \end{tabular*}\vspace{-7pt}
}

\newcommand{\resumeProjectHeading}[2]{
    \item
    \begin{tabular*}{0.97\textwidth}{l@{\extracolsep{\fill}}r}
      \small#1 & #2 \\
    \end{tabular*}\vspace{-7pt}
}

\newcommand{\resumeSubItem}[1]{\resumeItem{#1}\vspace{-4pt}}

\renewcommand\labelitemii{$\vcenter{\hbox{\tiny$\bullet$}}$}

\newcommand{\resumeSubHeadingListStart}{\begin{itemize}[leftmargin=0.15in, label={}]}
\newcommand{\resumeSubHeadingListEnd}{\end{itemize}}
\newcommand{\resumeItemListStart}{\begin{itemize}}
\newcommand{\resumeItemListEnd}{\end{itemize}\vspace{-5pt}}

% Role context under a \resumeSubheading. The heading table is top-aligned and
% pulls the next line up, so the blurb must start a new paragraph below it.
\newcommand{\resumeIntro}[1]{\par\vspace{3pt}{\small #1\par}}
.
% Author : Jake Gutierrez
% Based off of: https://github.com/sb2nov/resume
% License : MIT
% This file is not a document. Compile marti_ferran_gonzales_cv.tex or marti_ferran_gonzales_cover_letter.tex.

\ifdefined\cvsetupfont
  \usepackage{fontspec}
  \cvsetupfont
  \renewcommand{\familydefault}{\sfdefault}
\else
  \usepackage[T1]{fontenc}
  \usepackage[utf8]{inputenc}
\fi
\usepackage{latexsym}
\usepackage[empty]{fullpage}
\usepackage{titlesec}
\usepackage{marvosym}
\usepackage[usenames,dvipsnames]{color}
\usepackage{verbatim}
\usepackage{enumitem}
\usepackage[hidelinks]{hyperref}
\usepackage{fancyhdr}
\usepackage[english]{babel}
\usepackage{tabularx}
\ifdefined\cvsetupfont
\else
  \input{glyphtounicode}
\fi

\pagestyle{fancy}
\fancyhf{} % clear all header and footer fields
\fancyfoot{}
\renewcommand{\headrulewidth}{0pt}
\renewcommand{\footrulewidth}{0pt}

% Adjust margins
\addtolength{\oddsidemargin}{-0.5in}
\addtolength{\evensidemargin}{-0.5in}
\addtolength{\textwidth}{1in}
\addtolength{\topmargin}{-.5in}
\addtolength{\textheight}{1.0in}

\urlstyle{same}

\raggedbottom
\raggedright
\setlength{\tabcolsep}{0in}

% Sans faces here have no small caps, so the section title stays in the same font.
\ifdefined\cvsetupfont
  \titleformat{\section}{
    \vspace{-4pt}\raggedright\large
  }{}{0em}{\MakeUppercase}[\color{black}\titlerule \vspace{-5pt}]
\else
  \titleformat{\section}{
    \vspace{-4pt}\scshape\raggedright\large
  }{}{0em}{}[\color{black}\titlerule \vspace{-5pt}]
  % Ensure that the generated pdf is machine readable/ATS parsable
  \pdfgentounicode=1
\fi

%-------------------------
% Custom commands
\newcommand{\resumeItem}[1]{
  \item\small{
    {#1 \vspace{-2pt}}
  }
}

\newcommand{\resumeSubheading}[4]{
  \vspace{-2pt}\item
    \begin{tabular*}{0.97\textwidth}[t]{l@{\extracolsep{\fill}}r}
      \textbf{#1} & #2 \\
      \textit{\small#3} & \textit{\small #4} \\
    \end{tabular*}\vspace{-7pt}
}

\newcommand{\resumeSubSubheading}[2]{
    \item
    \begin{tabular*}{0.97\textwidth}{l@{\extracolsep{\fill}}r}
      \textit{\small#1} & \textit{\small #2} \\
    \end{tabular*}\vspace{-7pt}
}

\newcommand{\resumeProjectHeading}[2]{
    \item
    \begin{tabular*}{0.97\textwidth}{l@{\extracolsep{\fill}}r}
      \small#1 & #2 \\
    \end{tabular*}\vspace{-7pt}
}

\newcommand{\resumeSubItem}[1]{\resumeItem{#1}\vspace{-4pt}}

\renewcommand\labelitemii{$\vcenter{\hbox{\tiny$\bullet$}}$}

\newcommand{\resumeSubHeadingListStart}{\begin{itemize}[leftmargin=0.15in, label={}]}
\newcommand{\resumeSubHeadingListEnd}{\end{itemize}}
\newcommand{\resumeItemListStart}{\begin{itemize}}
\newcommand{\resumeItemListEnd}{\end{itemize}\vspace{-5pt}}

% Role context under a \resumeSubheading. The heading table is top-aligned and
% pulls the next line up, so the blurb must start a new paragraph below it.
\newcommand{\resumeIntro}[1]{\par\vspace{3pt}{\small #1\par}}
.
% Author : Jake Gutierrez
% Based off of: https://github.com/sb2nov/resume
% License : MIT
% This file is not a document. Compile marti_ferran_gonzales_cv.tex or marti_ferran_gonzales_cover_letter.tex.

\ifdefined\cvsetupfont
  \usepackage{fontspec}
  \cvsetupfont
  \renewcommand{\familydefault}{\sfdefault}
\else
  \usepackage[T1]{fontenc}
  \usepackage[utf8]{inputenc}
\fi
\usepackage{latexsym}
\usepackage[empty]{fullpage}
\usepackage{titlesec}
\usepackage{marvosym}
\usepackage[usenames,dvipsnames]{color}
\usepackage{verbatim}
\usepackage{enumitem}
\usepackage[hidelinks]{hyperref}
\usepackage{fancyhdr}
\usepackage[english]{babel}
\usepackage{tabularx}
\ifdefined\cvsetupfont
\else
  \input{glyphtounicode}
\fi

\pagestyle{fancy}
\fancyhf{} % clear all header and footer fields
\fancyfoot{}
\renewcommand{\headrulewidth}{0pt}
\renewcommand{\footrulewidth}{0pt}

% Adjust margins
\addtolength{\oddsidemargin}{-0.5in}
\addtolength{\evensidemargin}{-0.5in}
\addtolength{\textwidth}{1in}
\addtolength{\topmargin}{-.5in}
\addtolength{\textheight}{1.0in}

\urlstyle{same}

\raggedbottom
\raggedright
\setlength{\tabcolsep}{0in}

% Sans faces here have no small caps, so the section title stays in the same font.
\ifdefined\cvsetupfont
  \titleformat{\section}{
    \vspace{-4pt}\raggedright\large
  }{}{0em}{\MakeUppercase}[\color{black}\titlerule \vspace{-5pt}]
\else
  \titleformat{\section}{
    \vspace{-4pt}\scshape\raggedright\large
  }{}{0em}{}[\color{black}\titlerule \vspace{-5pt}]
  % Ensure that the generated pdf is machine readable/ATS parsable
  \pdfgentounicode=1
\fi

%-------------------------
% Custom commands
\newcommand{\resumeItem}[1]{
  \item\small{
    {#1 \vspace{-2pt}}
  }
}

\newcommand{\resumeSubheading}[4]{
  \vspace{-2pt}\item
    \begin{tabular*}{0.97\textwidth}[t]{l@{\extracolsep{\fill}}r}
      \textbf{#1} & #2 \\
      \textit{\small#3} & \textit{\small #4} \\
    \end{tabular*}\vspace{-7pt}
}

\newcommand{\resumeSubSubheading}[2]{
    \item
    \begin{tabular*}{0.97\textwidth}{l@{\extracolsep{\fill}}r}
      \textit{\small#1} & \textit{\small #2} \\
    \end{tabular*}\vspace{-7pt}
}

\newcommand{\resumeProjectHeading}[2]{
    \item
    \begin{tabular*}{0.97\textwidth}{l@{\extracolsep{\fill}}r}
      \small#1 & #2 \\
    \end{tabular*}\vspace{-7pt}
}

\newcommand{\resumeSubItem}[1]{\resumeItem{#1}\vspace{-4pt}}

\renewcommand\labelitemii{$\vcenter{\hbox{\tiny$\bullet$}}$}

\newcommand{\resumeSubHeadingListStart}{\begin{itemize}[leftmargin=0.15in, label={}]}
\newcommand{\resumeSubHeadingListEnd}{\end{itemize}}
\newcommand{\resumeItemListStart}{\begin{itemize}}
\newcommand{\resumeItemListEnd}{\end{itemize}\vspace{-5pt}}

% Role context under a \resumeSubheading. The heading table is top-aligned and
% pulls the next line up, so the blurb must start a new paragraph below it.
\newcommand{\resumeIntro}[1]{\par\vspace{3pt}{\small #1\par}}


% Role context under a \resumeSubheading. The heading table is top-aligned and
% pulls the next line up, so the blurb must start a new paragraph below it.
\newcommand{\resumeIntro}[1]{\par\vspace{3pt}{\small #1\par}}

\begin{document}

\begin{center}
    \textbf{\Huge MARTÍ FERRAN GONZALES} \\ \vspace{1pt}
    \small Cranfield, United Kingdom $|$ Vilafranca del Penedès, Spain \\ \vspace{1pt}
    \small +44 5673 224450 $|$
    \href{mailto:marti.ferran@icloud.com}{\underline{marti.ferran@icloud.com}} $|$
    \href{https://www.linkedin.com/in/mart\%C3\%AD-ferran-gonzales-080348229}{\underline{linkedin.com/in/martí-ferran-gonzales}}
\end{center}

\section{Personal Statement}
{\small Aerospace and Autonomous Systems Engineer specialising in Guidance, Navigation \& Control (GNC), flight dynamics, and real-time embedded systems. Experienced in developing closed-loop control algorithms, state estimation, and embedded software (C/C++, MATLAB/Simulink) validated through mathematical modelling, simulation, and dynamic testing on prototype platforms. Combines analytical work with a hands-on, collaborative approach to testing within multidisciplinary teams.  Focused on robust control architectures, high-fidelity dynamic modelling, and bridging theoretical simulation with physical prototype validation.}

\section{Key Achievements}
  \resumeItemListStart
    \resumeItem{Awarded Honours Distinction in Airport Design and Construction for achieving the highest grade in the module during the MSc in Aeronautical Engineering (Universitat Politècnica de Catalunya, 2026).}
    \resumeItem{Refactored the complete Vehicle Control Unit (VCU) software architecture in Simulink, optimising model execution, task scheduling, and communication reliability (CUPRA Racing, 2026).}
    \resumeItem{Built a complete multi-physics vehicle simulation model and developed a closed-loop yaw control algorithm, calibrating and verifying the system during prototype testing (CUPRA Racing, 2025).}
    \resumeItem{Coordinated the concept phase for the 2025 prototype vehicle, synthesising multidisciplinary ideas and aligning team efforts to converge on optimal vehicle targets (UPC ecoRacing, 2024).}
  \resumeItemListEnd

\section{Education}
  \resumeSubHeadingListStart
    \resumeSubheading
      {MSc in Autonomous Vehicle Dynamics and Control}{September 2026 -- Present}
      {Cranfield University (Cranfield, UK)}{}
      \resumeItemListStart
        \resumeItem{Key Modules: Guidance and Navigation for Autonomous Systems; Sensor Fusion; UAS Dynamics and Control; UAS Modelling and Simulation; Autonomous Vehicle Control Systems; AI for Autonomous Systems; Aerial Communication Systems}
      \resumeItemListEnd
    \resumeSubheading
      {MSc in Aeronautical Engineering}{September 2025 -- Present}
      {Universitat Politècnica de Catalunya (Terrassa, Spain)}{}
      \resumeItemListStart
        \resumeItem{Average Grade: 88.3\% (60 ECTS completed; 120 ECTS total)}
        \resumeItem{Key Modules: Aerodynamics, Flight and Orbital Mechanics; Air Transport and Navigation Systems; Aerospace Vehicles; RF and Communication Systems; Aircraft and Rocket Propulsion Systems; Airport Design and Construction (Honours Distinction, 2026)}
      \resumeItemListEnd
    \resumeSubheading
      {BSc in Aerospace Technologies Engineering}{September 2021 -- September 2025}
      {Universitat Politècnica de Catalunya (Terrassa, Spain)}{}
      \resumeItemListStart
        \resumeItem{Average Grade: 76.4\% (240 ECTS)}
        \resumeItem{Key Modules: Automatic Control; Flight Mechanics; Avionics; Electronic Circuits; Computational Engineering; Aerodynamics; Aerospace Structures; Aerospace Propulsion; Space Engineering; Aircraft Design}
        \resumeItem{Bachelor's Thesis: Study of yaw moment and stability control algorithm based on longitudinal torque distribution in electric vehicles --- Grade: 100\% (Nominated for Honours Distinction)}
      \resumeItemListEnd
  \resumeSubHeadingListEnd

\section{Career History}
  \resumeSubHeadingListStart
    \resumeSubheading
      {Electronics and Embedded Software Engineering Intern}{January 2025 -- February 2026}
      {CUPRA Racing (Abrera, Spain)}{}
      \resumeIntro{Motorsport division of CUPRA, responsible for the development of high-performance racing vehicles, experimental prototype platforms, and advanced engineering R\&D projects.}
      \resumeItemListStart
        \resumeItem{Built vehicle multi-physics simulation models in MATLAB/Simulink for the Touring Car Racing (TCR) prototype and its hybrid version, integrating non-linear tyre dynamics, suspension dynamics, and hybrid powertrain systems (ICE and electric motor).}
        \resumeItem{Formulated and simulated closed-loop yaw stability control algorithms based on driving and braking torque distribution for TCR and Raval 4WD prototypes; tuned and assessed control laws during track testing to guarantee dynamic stability and maximum performance.}
        \resumeItem{Refactored the complete Vehicle Control Unit (VCU) software architecture in Simulink for the TCR hybrid platform, optimising model execution, task scheduling, and real-time communication protocols.}
        \resumeItem{Modelled, simulated, and sized an electric powertrain cooling system in MATLAB/Simulink, achieving high model fidelity correlated against empirical sensor data.}
        \resumeItem{Configured data-acquisition systems and sensor instrumentation, collaborating with trackside engineers and test drivers to monitor dynamic parameters and diagnose live prototype performance.}
      \resumeItemListEnd

    \resumeSubheading
      {Vehicle Dynamics \& Control Engineer}{November 2021 -- September 2024}
      {UPC ecoRacing (Formula Student) (Terrassa, Spain)}{}
      \resumeIntro{Spain's pioneering and most awarded Formula Student engineering team, designing and manufacturing high-performance 4WD electric and autonomous (manual and driverless) formula prototypes.}
      \resumeItemListStart
        \resumeItem{Engineered a complete closed-loop torque vectoring algorithm for a 4WD electric prototype with four independent in-wheel motors, maximising cornering yaw response through torque distribution and conducting live-session parameter tuning.}
        \resumeItem{Led the concept design phase for the 2025 prototype vehicle, coordinating multidisciplinary departments across vehicle dynamics, chassis, powertrain, and aerodynamics to converge on balanced vehicle targets.}
        \resumeItem{Created non-linear vehicle dynamics simulation models in MATLAB and Python, integrating semi-empirical tyre models and parameter identification algorithms from test data to improve simulation accuracy and support concept and setup optimisation.}
        \resumeItem{Served as Test and Validation Lead, planning and executing testing sessions to verify electronic control algorithms, streamline energy management, and validate mechanical setup and vehicle dynamics performance.}
        \resumeItem{Modelled suspension components in CAD, assessed load cases, validated structural integrity via FEA, and conducted on-track testing using strain-gauge instrumentation to correlate dynamic stress distributions as well as the boundary conditions used.}
      \resumeItemListEnd

    \resumeSubheading
      {Head of Logistics}{September 2023 -- September 2024}
      {UPC ecoRacing (Formula Student) (Terrassa, Spain)}{}
      \resumeItemListStart
        \resumeItem{Liaised with technical sub-teams and competition scrutineers to ensure seamless equipment readiness, trackside operations, and schedule adherence across international Formula Student events.}
        \resumeItem{Managed testing schedules, operational timelines, competition travel logistics, and inventory allocation for a 40+ member engineering team.}
      \resumeItemListEnd
  \resumeSubHeadingListEnd

\section{Technical Projects}
  \resumeSubHeadingListStart
    \resumeSubheading
      {Vehicle Dynamics \& Telemetry Analysis Suite}{July 2026 -- Present}
      {System Architecture \& AI-Assisted Development}{}
      \resumeItemListStart
        \resumeItem{Definined the system architecture and engineering logic for a modular desktop software suite focused on telemetry data analysis and suspension kinematics evaluation.}
        \resumeItem{Coordinated development through structured AI-agent workflows to implement data import pipelines, session comparison tools, custom mathematical channels, and visualisation layouts with minimal manual coding effort.}
      \resumeItemListEnd
    \resumeSubheading
      {Nutria (Desktop Clinical Management Platform)}{September 2024 -- August 2026}
      {Python, Product Design \& AI Workflows}{}
      \resumeItemListStart
        \resumeItem{Programmed the initial functional prototypes in Python and subsequently led the product architecture and system redesign from concept to a private pilot release.}
        \resumeItem{Designed and orchestrated AI-driven engineering workflows, structuring system requirements, modular architecture, and automated End-to-End testing to implement features with minimal human intervention and minimum iterations.}
      \resumeItemListEnd
  \resumeSubHeadingListEnd

\section{Skills}
 \begin{itemize}[leftmargin=0.15in, label={}]
    \small{\item{
     \textbf{Programming:} MATLAB, Simulink, C, C++, Python, LaTeX. \\
     \textbf{Engineering, Simulation \& Telemetry:} SolidWorks, CATIA, FreeCAD, MSC Nastran, ANSA, MoTeC, WinTAX4, Sysma. \\
     \textbf{Agentic-AI \& Developer Tools:} Cursor, Antigravity, Claude Code, Git, GitHub, VS Code. \\
     \textbf{Languages:} Catalan (Native), Spanish (Native), English (C2 Proficient), German (A1 Basic).
    }}
 \end{itemize}

\section{Training}
 \begin{itemize}[leftmargin=0.15in, label={}]
    \small{\item{
     \textbf{Applied Vehicle Dynamics Seminar} --- Optimum G (Stuttgart, Germany $|$ November 2024, 48 hours): Intensive training on non-linear tyre modelling, steady-state and transient dynamics, load transfer, kinematics and compliance (K\&C), and track simulation validation. \\
     \vspace{1pt}
     \textbf{Data-Driven Performance Engineering Seminar} --- Optimum G (Stuttgart, Germany $|$ November 2024, 48 hours): Practical seminar on data-acquisition hardware/software, sensor analysis, track testing methodologies, and data-driven engineering decision-making.
    }}
 \end{itemize}

\end{document}
